\chapter*{Kata Pengantar}
\addcontentsline{toc}{chapter}{KATA PENGANTAR}

Puji dan syukur penulis panjatkan kepada Tuhan Yang Maha Esa atas berkat dan rahmat-Nya sehingga laporan tugas akhir yang berjudul ``\thetitle'' dapat diselesaikan dalam rangka memenuhi syarat kelulusan tingkat sarjana. Dalam pengerjaan tugas akhir ini, penulis memperoleh banyak dukungan dari berbagai pihak. Untuk itu, penulis mengucapkan terima kasih kepada:

\begin{enumerate}
  \item Dr.\ Agung Dewandaru, S.T., M.Sc.\ dan I Wayan Gunada, S.H., M.H.\ selaku dosen pembimbing atas bimbingan, arahan, dan kesabarannya selama pengerjaan tugas akhir.
  \item Seluruh dosen Program Studi Teknik Informatika ITB atas ilmu pengetahuan yang telah diberikan.
  \item Rekan-rekan satu tim \textit{capstone} ``OMNI-RAG'' atas kerja sama dan kolaborasinya selama pengembangan sistem.
  \item Kedua orang tua dan keluarga penulis atas dukungan moral, motivasi, dan doa yang tiada henti.
  \item Max Kusnadi selaku manajer dan Nitin Pillai selaku pemimpin tim di Xendit atas izin cuti selama dua minggu sehingga penulis dapat berfokus menyelesaikan tugas akhir ini.
  \item Malik Irawan dan rekan-rekan asisten Laboratorium Teknik Produksi FTMD yang telah menemani penulis mengerjakan tugas akhir hingga larut malam.
  \item Seluruh pihak lain yang tidak dapat penulis sebutkan satu per satu yang telah membantu dalam proses pengerjaan tugas akhir ini.
\end{enumerate}

Penulis menyadari bahwa laporan tugas akhir ini masih memiliki kekurangan. Oleh karena itu, kritik dan saran yang membangun sangat penulis harapkan. Semoga laporan ini dapat bermanfaat bagi pembaca dan pengembangan penelitian selanjutnya.

\begin{flushright}
  \vspace{0.5cm}
  Bandung, \tanggalpengesahan

  \vspace{1.5cm}

  Dewantoro Triatmojo \\
  NIM 13522011
\end{flushright}
